% ============================================================
% V. DISCUSIÓN
% ============================================================

\section{Discusión}

Los resultados muestran que la Transformada Rápida de Fourier constituye una herramienta útil para analizar señales de audio y observar diferencias entre diferentes fuentes sonoras. La frecuencia dominante fue, de las tres características calculadas, la que presentó mayor utilidad para distinguir las muestras analizadas de forma cualitativa: en los tres grupos (animales, instrumentos, voces) las frecuencias dominantes estuvieron suficientemente separadas para observar diferencias claras entre las señales de cada grupo (Figs.~\ref{fig:animales_comparacion}, \ref{fig:instrumentos_comparacion}, \ref{fig:personas_comparacion}).

En el caso de los instrumentos, el contrabajo presentó la frecuencia dominante más baja de las ocho señales, como es de esperar de un instrumento grave, pero su valor debe leerse con cautela: la nota solo suena en los últimos $0.5$ s de la ventana de $3$ s analizada (Fig.~\ref{fig:contrabajo}), por lo que el resultado describe un ataque corto, no una nota sostenida. El piano, al reproducir una frase de varias notas en vez de una sola, obtuvo la frecuencia dominante más alta de los tres instrumentos —por encima de la flauta—, lo que confirma que ``frecuencia dominante'' identifica el pico de mayor energía dentro de la ventana analizada, y no la nota fundamental de un evento sonoro aislado; comparar frecuencias dominantes entre grabaciones con contenido musical distinto (una nota vs. una frase) es, por tanto, una comparación de la ventana capturada, no del registro característico del instrumento.

\subsection{Una debilidad de unidades en el puntaje de separabilidad}
\label{sec:discusion-unidades}

Los siete pares reales de la Tabla~\ref{tab:separabilidad} —tres de animales, tres de instrumentos, uno de voces— resultaron \emph{todos} claramente separables, con puntajes entre $442$ y $7788$. A primera vista esto podría leerse como una confirmación general de que la frecuencia dominante distingue bien cualquier par de fuentes sonoras. Sin embargo, la uniformidad del resultado revela más bien una debilidad de la Ec.~\ref{eq:separabilidad}: al dividir una diferencia de frecuencia dominante (en Hz, típicamente cientos o miles) por una suma de desviaciones estándar de amplitud normalizada por \texttt{librosa} (típicamente $\sigma \in [0.01, 0.3]$), el denominador es sistemáticamente pequeño para cualquier grabación de audio real, de modo que prácticamente cualquier par con frecuencias dominantes distintas por más de un puñado de Hz supera con holgura el umbral de $3$. Dicho de otro modo: una vez fijado el pipeline de normalización de amplitud (como aquí, para las ocho señales), el puntaje deja de discriminar entre pares ``muy separables'' y ``apenas separables'' — todos caen en la misma categoría.

Esto es consistente con lo observado al recalcular la Tabla~\ref{tab:separabilidad} en una revisión anterior de este informe, cuando los mismos tres pares de animales, evaluados sobre valores de referencia preliminares con desviaciones estándar de otra escala ($10^{3}$ en vez de $10^{-1}$), daban puntajes de $0.11$--$0.33$ (``solapados''). Es decir, el mismo formulismo produce, según la escala del denominador, o bien un ``todo solapado'' o bien un ``todo separable'' — ninguno de los dos extremos informativo por sí solo. Una versión más robusta del puntaje debería normalizar ambos términos a magnitudes comparables: por ejemplo, dividiendo $\Delta f_d$ por el ancho de banda de Nyquist ($f_s/2$) en lugar de por $\sigma_a+\sigma_b$, o sustituyendo el denominador por la dispersión de la propia frecuencia dominante estimada sobre varias subventanas de cada grabación, de forma que numerador y denominador compartan unidades de frecuencia.

Por otra parte, la media de las ocho señales resultó cercana a cero en todos los casos —incluidas las dos voces sintetizadas, con medias de $0.0196$ y $0.0265$—, como es de esperar para audio sin componente de continua relevante. Ni la media ni la desviación estándar por sí solas resultaron adecuadas para diferenciar las fuentes sonoras: solo adquieren utilidad como parte del denominador de un puntaje de separabilidad correctamente escalado, no como características independientes.

Debe considerarse además que la magnitud de la FFT depende de factores como la amplitud de grabación, la duración de la señal y las condiciones de captura, por lo que no debe utilizarse directamente como criterio de identificación sin una normalización previa consistente entre grabaciones —como la que aquí garantiza \texttt{librosa.load} al cargar las ocho señales por igual.

Finalmente, dos de las tres categorías (animales, instrumentos) usan audio real, pero solo una grabación por señal, y la categoría de voces usa habla sintetizada en lugar de grabaciones humanas, lo que limita cualquier generalización sobre variabilidad de voz real. El número de muestras impide, en los tres casos, determinar si las diferencias observadas son estadísticamente significativas más allá de estos ejemplos puntuales. Para desarrollar un sistema de identificación más robusto sería necesario ampliar \texttt{audio\_samples/} con múltiples grabaciones reales por categoría —incluyendo voces humanas reales—, corregir la Ec.~\ref{eq:separabilidad} como se propone arriba, y considerar características adicionales a la frecuencia dominante y la dispersión global, como la distribución de armónicos o la energía por bandas.
